% GENERATED FILE. Edit canonical lessons, then run npm run generate:books.
% canonical-source-hash: fnv1a64:15bc9b3e2bd75c5c
% canonical-lessons: TE-C101-paina, TE-C101-kinda, TE-C101-lopala, TE-C101-bayata, TE-C101-venaka

\chapter{Above, Below, Inside, Outside, Behind}
\label{ch:te-above-below-inside-outside-behind}
\hlchaptermodality{101}

\begin{chapteropening}
\textbf{By the end of this chapter:} \emph{I can say above, below, inside, outside and behind.}

Complete the last lesson of chapter 101: i can say above, below, inside, outside and behind.
\end{chapteropening}

\section[paina]{\te{పైన} (paina) — above, on top}

\label{lesson:TE-C101-paina}



\begin{quote}
Before the new one: say the Telugu for near, then the Telugu for far; distance.
\end{quote}

\subsection*{You'll want to know: \te{పైన}}
\textbf{\te{పైన}} — \emph{paina} — \textquotedblleft{}above, on top\textquotedblright{}.

\begin{grammarlens}[title={What you've built}]
One more word for this chapter.
\end{grammarlens}

\subsection*{Your turn}
\begin{itemize}
\raggedright
  \item \emph{Say it:} \emph{paina}
  \item \emph{Say it:} \emph{paina}, once more
  \item \emph{Say it:} say \emph{paina}
  \item \emph{Recall:} say the Telugu for near, then the Telugu for far; distance, then say \emph{paina} again
\end{itemize}

\begin{tcolorbox}[breakable,colback=teal!4,colframe=teal!35!black,title={Before you move on}]
What does \te{పైన} mean? (\textquotedblleft{}Above, on top\textquotedblright{}.) Say it once more.
\end{tcolorbox}
\section[kinda]{\te{కింద} (kinda) — below, underneath}

\label{lesson:TE-C101-kinda}



\begin{quote}
Before the new one: say the Telugu for far; distance, then the Telugu for above.
\end{quote}

\subsection*{You'll want to know: \te{కింద}}
\textbf{\te{కింద}} — \emph{kinda} — \textquotedblleft{}below, underneath\textquotedblright{}.

\begin{grammarlens}[title={What you've built}]
One more word for this chapter.
\end{grammarlens}

\subsection*{Your turn}
\begin{itemize}
\raggedright
  \item \emph{Say it:} \emph{kinda}
  \item \emph{Say it:} \emph{kinda}, once more
  \item \emph{Say it:} say \emph{kinda}
  \item \emph{Recall:} say the Telugu for far; distance, then the Telugu for above, then say \emph{kinda} again
\end{itemize}

\begin{tcolorbox}[breakable,colback=teal!4,colframe=teal!35!black,title={Before you move on}]
What does \te{కింద} mean? (\textquotedblleft{}Below, underneath\textquotedblright{}.) Say it once more.
\end{tcolorbox}
\section[lōpala]{\te{లోపల} (lōpala) — inside}

\label{lesson:TE-C101-lopala}



\begin{quote}
Before the new one: say the Telugu for above, then the Telugu for below.
\end{quote}

\subsection*{You'll want to know: \te{లోపల}}
\textbf{\te{లోపల}} — \emph{lōpala} — \textquotedblleft{}inside\textquotedblright{}.

\begin{grammarlens}[title={What you've built}]
One more word for this chapter.
\end{grammarlens}

\subsection*{Your turn}
\begin{itemize}
\raggedright
  \item \emph{Say it:} \emph{lōpala}
  \item \emph{Say it:} \emph{lōpala}, once more
  \item \emph{Say it:} say \emph{lōpala}
  \item \emph{Recall:} say the Telugu for above, then the Telugu for below, then say \emph{lōpala} again
\end{itemize}

\begin{tcolorbox}[breakable,colback=teal!4,colframe=teal!35!black,title={Before you move on}]
What does \te{లోపల} mean? (\textquotedblleft{}Inside\textquotedblright{}.) Say it once more.
\end{tcolorbox}
\section[bayaṭa]{\te{బయట} (bayaṭa) — outside}

\label{lesson:TE-C101-bayata}



\begin{quote}
Before the new one: say the Telugu for below, then the Telugu for inside.
\end{quote}

\subsection*{You'll want to know: \te{బయట}}
\textbf{\te{బయట}} — \emph{bayaṭa} — \textquotedblleft{}outside\textquotedblright{}.

\begin{grammarlens}[title={What you've built}]
One more word for this chapter.
\end{grammarlens}

\subsection*{Your turn}
\begin{itemize}
\raggedright
  \item \emph{Say it:} \emph{bayaṭa}
  \item \emph{Say it:} \emph{bayaṭa}, once more
  \item \emph{Say it:} say \emph{bayaṭa}
  \item \emph{Recall:} say the Telugu for below, then the Telugu for inside, then say \emph{bayaṭa} again
\end{itemize}

\begin{tcolorbox}[breakable,colback=teal!4,colframe=teal!35!black,title={Before you move on}]
What does \te{బయట} mean? (\textquotedblleft{}Outside\textquotedblright{}.) Say it once more.
\end{tcolorbox}
\section[venaka]{\te{వెనక} (venaka) — behind}

\label{lesson:TE-C101-venaka}



\begin{quote}
Before the new one: say the Telugu for inside, then the Telugu for outside.
\end{quote}

\subsection*{You'll want to know: \te{వెనక}}
\textbf{\te{వెనక}} — \emph{venaka} — \textquotedblleft{}behind\textquotedblright{}.

\begin{grammarlens}[title={What you've built}]
One more word for this chapter.
\end{grammarlens}

\subsection*{Your turn}
\begin{itemize}
\raggedright
  \item \emph{Say it:} \emph{venaka}
  \item \emph{Say it:} \emph{venaka}, once more
  \item \emph{Say it:} say \emph{venaka}
  \item \emph{Recall:} say the Telugu for inside, then the Telugu for outside, then say \emph{venaka} again
\end{itemize}

\begin{tcolorbox}[breakable,colback=teal!4,colframe=teal!35!black,title={Before you move on}]
What does \te{వెనక} mean? (\textquotedblleft{}Behind\textquotedblright{}.) Say it once more.
\end{tcolorbox}
